\section{Modèles de copules pour le risque de défaut corrélé}

Nous avons vu section \ref{subsubsec:modele_copule} qu'il est possible --- grâce au théorème de Sklar ---
de reconstituer les dépendances entre les marginales via des copules afin d'obtenir la loi conjointe 
des temps de défauts. Bien que la copule exacte soit inconnue il est possible d'utiliser des modèles
paramétriques afin d'approximer cette structure de dépendances. Cette section aborde différentes
classes de copules adaptées aux modèles de défauts corrélés.

\subsection{La copule gaussienne à un facteur}\label{subsec:copule_gauss}

La copule gaussienne constitue l'approche historiquement dominante pour la tarification
des tranches de CDO depuis les travaux de Li (2000). Considérons les marginales 
$F_i(t) = \mathbb{P}^*(\tau_i \leq t)$, la copule gaussienne suppose que les variables
$\Phi^{-1}(F_i(\tau_i))$ suivent conjointement une loi normale multivariée de matrice de
corrélation $\Sigma$, ce qui s'écrit
\begin{equation}
    C^\text{Gauss}_\Sigma(u_1, \ldots, u_n) = \Phi_\Sigma\bigl(\Phi^{-1}(u_1), \ldots, \Phi^{-1}(u_n)\bigr),
\end{equation}
où $\Phi$ désigne la fonction de répartition de la loi normale centrée réduite et $\Phi_\Sigma$ celle de
$\mathcal{N}(0, \Sigma)$. Sous l'hypothèse de portefeuille homogène, on impose une corrélation uniforme
$\rho$ entre toutes les paires d'entités, réduisant le problème à l'estimation d'un seul paramètre.
La structure de dépendance se représente alors par des variables latentes corrélées sous un modèle à un facteur commun $M \sim \mathcal{N}(0,1)$ :
\begin{equation}
    X_i = \sqrt{\rho}\, M + \sqrt{1-\rho}\, \varepsilon_i,
    \label{eq:facteur_gauss}
\end{equation}
où les $\varepsilon_i \sim \mathcal{N}(0,1)$ sont indépendants entre eux et indépendants de $M$.
Nous pouvons remarquer que $X_i$ est une combinaison linéaire de gaussiennes et donc est elle-même une
gaussienne qui, de plus, est centrée réduite. Nous retrouvons également que pour tout $i\neq j$, $\text{Cov}(X_i,X_j)=\rho$.

Le défaut de l'entité $i$ avant la date $t$ est alors défini par l'évènement $\{X_i \leq c_i(t)\}$ où $c_i(t)$
est un seuil à déterminer en faisant coïncider la probabilité de cet évènement avec celle donnée par la marginale
correspondante :
\begin{equation}
    \mathbb{P}^\ast (X_i \leq c_i(t)) = F_i(t) .
\end{equation}
Une solution analytique est ici possible pour l'équation et nous obtenons $c_i(t) = \Phi^{-1}(F_i(t))$ et l'évènement
de défaut de l'entité $i$ avant la date $t$ est donné par $X_i \leq \Phi^{-1}(F_i(t))$.

La copule gaussienne ne présente aucune dépendance de queue : lorsque $\rho$ est fixé,
la probabilité que deux entités fassent défaut simultanément lors d'un événement extrême tend
vers zéro, ce qui constitue une limitation majeure face aux données observées en période de stress.

\subsection{La copule de Student}

Une extension naturelle de la copule gaussienne est la copule de Student, qui s'écrit
\begin{equation}
    C^t_{\Sigma,\nu}(u_1, \ldots, u_n) = t_{\Sigma,\nu}\bigl(t_\nu^{-1}(u_1), \ldots, t_\nu^{-1}(u_n)\bigr),
\end{equation}
où $t_\nu$ désigne la fonction de répartition de la loi de Student univariée à $\nu$ degrés de
liberté, et $t_{\Sigma,\nu}$ celle de la loi de Student multivariée. Le paramètre $\nu$ contrôle
l'épaisseur des queues de distribution : plus $\nu$ est faible, plus les défauts simultanés
extrêmes sont favorisés. Lorsque $\nu \to \infty$, la copule de Student converge vers la copule
gaussienne. Sa simulation repose sur la représentation stochastique suivante : on génère un
vecteur gaussien $Z = (Z_1, \ldots, Z_n)$ de matrice de corrélation $\Sigma(\rho)$ via le même
facteur commun~\eqref{eq:facteur_gauss}, puis on pose
\begin{equation}
    T_i = \frac{Z_i}{\sqrt{S/\nu}} \sim \mathcal{T}(\nu), \quad U_i = t_\nu(T_i),
\end{equation}
où $S \sim \chi^2(\nu)$ est une variable commune indépendante des $Z_i$ et $\mathcal{T}(\nu)$ désigne la
loi de Student à $\nu$ degrés de liberté. Le vecteur $(U_i)$ suit
alors une copule de Student de paramètres $(\rho, \nu)$.

Comme précédemment le défaut de l'entité $i$ avant la date $t$ est défini comme l'évènement $\{T_i \leq c_i(t) \}$
avec cette fois $c_i(t) = t_\nu^{-1}(F_i(t))$.

\subsection{La copule t-gaussienne à un facteur (Hull-White)}

La copule de Student attribue des queues épaisses à la fois au facteur systémique et aux
chocs idiosyncratiques. Hull et White (2004) proposent une variante dans laquelle seul le
facteur systémique présente des queues épaisses, les termes idiosyncratiques restant gaussiens.
Le modèle à un facteur s'écrit
\begin{equation}
    X_i = \sqrt{\rho}\, \tilde{M} + \sqrt{1-\rho}\, \varepsilon_i,
\end{equation}
où $\tilde{M} \sim \mathcal{T}(\nu)$ est un facteur commun de Student à $\nu$ degrés de liberté et les
$\varepsilon_i \sim \mathcal{N}(0,1)$ sont des termes idiosyncratiques gaussiens indépendants.
Ce choix se justifie économiquement : les chocs de marché systémiques (crises de liquidité,
effondrements sectoriels) ont une nature différente des défauts individuels, et leur modélisation
par une loi à queues épaisses permet de générer des épisodes de défauts simultanés sans pour
autant modifier le comportement de chaque entité prise isolément. 

Une extension naturelle de ce modèle est la copule t-t (ou copule double t) pour laquelle
$\tilde{M} \sim \mathcal{T}(\nu_1)$ est un facteur commun de Student à $\nu_1$ degrés de liberté 
et les $\varepsilon_i \sim \mathcal{T}(\nu_2)$ sont des termes idiosyncratiques indépendants de Student.
Cependant les tentatives de calibrations sont coûteuses et montrent dans notre cas d'étude que $\nu_2$
devient très grand systématiquement et donc les $\varepsilon_i$ tendent vers des lois gaussiennes.
Pour cette raison nous excluons ce modèle bien qu'il soit présent dans la littérature et nous
resterons sur la copule t-gaussienne à un facteur.

La variable $X_i$ n'a pas de loi connue explicite, mais conditionnellement
au facteur $\tilde{M} = m$, les termes idiosyncratiques restant gaussiens, on obtient
\begin{equation}
    X_i \mid \tilde{M} = m \;\sim\; \mathcal{N}\!\left(\sqrt{\rho}\, m,\; 1-\rho\right).
\end{equation}
Ainsi l'évènement de défaut de l'entité $i$ avant $t$ conditionnellement au facteur commun $\tilde{M} = m$
est donné par $\{ X_i \leq c_i(t) \mid \tilde{M} = m\}$ et en remarquant que 
$Z_i = \frac{X_i - \sqrt{\rho}\, m,}{\sqrt{1-\rho}}$ est centrée réduite la probabilité de défaut conditionnelle
s'écrit alors
\begin{equation}
    \mathbb{P}^\ast\bigl(\tau_i \leq t \mid \tilde{M} = m\bigr)
    =\mathbb{P}^\ast\bigl(X_i \leq c_i(t) \mid \tilde{M} = m\bigr)
    = \Phi\!\left(\frac{c_i(t) - \sqrt{\rho}\, m}{\sqrt{1-\rho}}\right) .
\end{equation}
Pour retrouver la probabilité de défaut marginale il suffit de prendre l'espérance risque neutre sur $\tilde{M}$
cependant il est à présent impossible de trouver une solution analytique pour $c_i(t)$. Ainsi le seuil est 
obtenu par inversion numérique de la fonction de répartition marginale de $X_i$ : 
\begin{equation}
    F_{X_i}(x) = \mathbb{E}^{\mathbb{P}^\ast}_{\tilde{M}}\left[\Phi\left(\frac{x - \sqrt{\rho}\,\tilde{M}}{\sqrt{1-\rho}}\right)\right] .
\end{equation}
Cette inversion est réalisée pour chaque date de paiement par la méthode de Newton-Raphson,
comme détaillé en section~\ref{sec:methodes_numeriques}.

\subsection{Le modèle de mélange de corrélations (BGL)}

Burtschell, Gregory et Laurent (2009) proposent un modèle dans lequel la corrélation
$\rho_B$ n'est pas fixe mais tirée d'un mélange discret à deux régimes. L'idée est de représenter
explicitement une bimodalité entre un régime normal, où les défauts sont relativement indépendants,
et un régime de crise, où ils sont fortement corrélés. Le facteur de risque systémique s'écrit
\begin{equation}
    \rho_B = \begin{cases} \rho_1 & \text{avec probabilité } q, \\ \rho_2 & \text{avec probabilité } 1-q, \end{cases}
\end{equation}
avec $0 < \rho_1 < \rho_2 < 1$ et $q \in (0,1)$. Conditionnellement au régime, la structure
de dépendance reste gaussienne à un facteur, de sorte que
\begin{equation}
    X_i = \sqrt{\rho_B}\, M + \sqrt{1-\rho_B}\, \varepsilon_i,
\end{equation}
où $M, \varepsilon_i \sim \mathcal{N}(0,1)$ indépendants. Ce modèle à trois paramètres $(\rho_1, \rho_2, q)$
autorise une flexibilité considérablement accrue pour reproduire simultanément les tranches
equity et les tranches senior, qui requièrent des niveaux de dépendance opposés dans le cadre
gaussien à un facteur. Le seuil définissant l'évènement de défaut est ici le même que pour le
modèle à un facteur gaussien décrit section \ref{subsec:copule_gauss}, il en sera de même pour le
modèle à corrélation stochastique décrit \ref{subsec:copule_cor_sto}.

Sous l'approximation de grand portefeuille homogène (LHP, voir
section~\ref{sec:methodes_numeriques}), la perte espérée de tranche se décompose comme
\begin{equation}
    \mathrm{ETL}_\text{BGL}(t) = q \cdot \mathrm{ETL}_\text{Vasicek}(\rho_1, t)
    + (1-q) \cdot \mathrm{ETL}_\text{Vasicek}(\rho_2, t),
\end{equation}
ce qui ramène chaque évaluation à deux appels à la formule analytique du modèle de Vasicek,
sans simulation Monte Carlo.

\subsection{Le modèle de corrélation stochastique (Vasicek dynamique)}\label{subsec:copule_cor_sto}

Une alternative continue au mélange discret est le modèle dit de corrélation stochastique,
dans lequel $\rho$ est elle-même une variable aléatoire suivant une loi Bêta de paramètres
$(\bar{\rho}\,\kappa, (1-\bar{\rho})\kappa)$, avec $\bar{\rho} \in (0,1)$ et $\kappa > 0$. Ce choix
de loi est motivé par le fait que la distribution Bêta est de support $[0,1]$ et
admet des formes unimodales, bimodales ou uniformes selon ses paramètres. On rappelle que pour
$\rho \sim \text{Bêta}(a, b)$ avec $a = \bar{\rho}\,\kappa$ et $b = (1-\bar{\rho})\kappa$,
\begin{equation}
    \mathbb{E}^{\mathbb{P}^\ast}[\rho] = \bar{\rho}, \qquad \mathrm{Var}(\rho) = \frac{\bar{\rho}(1-\bar{\rho})}{\kappa+1},
\end{equation}
de sorte que $\kappa$ contrôle la concentration de la distribution autour de sa moyenne : lorsque
$\kappa \to \infty$, la corrélation tend presque sûrement vers $\bar{\rho}$ et le modèle converge
vers le modèle de Vasicek déterministe. Cette loi Beta est la distribution stationnaire du processus 
de \emph{mean-reversion} contraint à $[0,1]$, appelé processus de diffusion de Jacobi (processus de 
Ornstein-Uhlenbeck avec bruit multiplicatif) :
\begin{equation}
    d\rho_t = \kappa(\bar\rho - \rho_t)\,dt + \sigma_\rho\sqrt{\rho_t(1-\rho_t)}\,dW_t
\end{equation}
Conditionnellement à $\rho$, la structure de dépendance reste le modèle gaussien à un facteur.
La perte espérée de tranche s'obtient alors par intégration sur la loi de $\rho$ :
\begin{equation}
    \mathrm{ETL}_\text{SR}(t) = \int_0^1 \mathrm{ETL}_\text{Vasicek}(\rho, t)\, f_\text{Bêta}(\rho)\, d\rho,
\end{equation}
où $f_\text{Bêta}$ est la densité de la loi $\text{Bêta}(\bar{\rho}\,\kappa, (1-\bar{\rho})\kappa)$.
Cette intégrale unidimensionnelle est évaluée par quadrature de Gauss-Legendre, chaque nœud
faisant lui-même appel à la formule analytique de Vasicek. Le modèle à deux paramètres $(\bar{\rho}, \kappa)$
offre ainsi une représentation continue de l'incertitude sur la corrélation, là où le modèle BGL
procède par mélange discret.
